\subsubsection{覆盖内容、原理与求解步骤}

队长的模型链从守恒关系开始，而不是从求解器开始。机械系统用牛顿第二定律、
刚体平衡、重力、摩擦和弹簧—阻尼关系建立状态方程；流体问题用质量守恒和
Bernoulli 关系闭合；电路问题用 Ohm 定律和 Kirchhoff 定律组织节点或回路方程；
气体过程核对理想气体状态关系。传热采用 Newton 冷却或热传导方程，波动、扩散
及多孔介质渗流按偏微分方程处理。人口过程用 Logistic 方程和捕食—被捕食模型，
药物过程使用房室代谢方程，传播过程比较 SIR 与 SIS 的恢复机制。

完整步骤固定为：
\begin{enumerate}
  \item 画出对象、边界和能量/质量/动量交换，确定系统尺度与正方向；
  \item 从守恒律或受力关系推导状态方程，逐项给出单位和符号；
  \item 写明初始条件、边界条件、连续性条件、控制量及可行域；
  \item 先进行极端情形、量纲和简化模型的解析检查；
  \item 对 ODE 比较显式 Euler 与四阶 Runge--Kutta，对 PDE 按网格选择有限差分或有限元；
  \item 进行步长/网格收敛、守恒残差、回代误差和参数扰动检验；
  \item 把经检验的状态量封装为组员1的目标/约束或组员2的可解释特征。
\end{enumerate}

Euler 法实现简单但全局误差阶数低，只用于基线和步长敏感性对照；RK4 在光滑 ODE
上兼顾精度与实现成本，但对刚性系统不应硬套。有限差分适合规则区域和清晰的
差分格式，可直接检查截断误差与 CFL 条件；有限元更适合复杂几何和变系数边界，
代价是网格、形函数、装配和边界处理更重。解析解存在时先用解析解作单元测试，
不存在时采用网格加密、守恒量和独立实现交叉验证。

